\SourceLeaf{292}%
a triangle side \(r\), its azimuth \(T\) at the starting point, and the latitude \(S\) of that point, derive the azimuth \(T'\pm180^0\) of the side at the other endpoint, its latitude \(S'\), and the difference in longitude \(\lambda\) between the two points. This amounts to solving a spherical triangle. The problem deserves a place here only because the usual formulas need some rearrangement if, in view of the remark in Art. 10, we want results as accurate as the given \(r\) without using logarithms with more digits. So that we can choose among the methods of solution according to the needs of each case, I first give those based on the familiar elementary formulas of spherical trigonometry.

First method
\[
\begin{aligned}
\operatorname{tang}s&=\cos T\operatorname{tang}r\\
\operatorname{tang}\lambda&=\frac{\operatorname{tang}T\sin s}{\cos(S-s)}\\
\operatorname{tang}S'&=\cos\lambda\operatorname{tang}(S-s)\\
\sin T'&=\frac{\sin T\cos S}{\cos S'}
\end{aligned}
\]
\indent Second method
\[
\begin{aligned}
\operatorname{tang}R&=\frac{\operatorname{tang}S}{\cos T}\\
\operatorname{tang}T'&=\frac{\operatorname{tang}T\cos R}{\cos(R-r)}\\
\operatorname{tang}S'&=\cos T'\operatorname{tang}(R-r)\\
\sin\lambda&=\frac{\sin r\sin T}{\cos S'}=\frac{\sin r\sin T'}{\cos S}
\end{aligned}
\]
\indent Third method
\[
\begin{aligned}
\sin(45^0+\tfrac{1}{2}S')\sin\tfrac{1}{2}(T'+\lambda)&=\sin(45^0+\tfrac{1}{2}(S+r))\sin\tfrac{1}{2}T\\
\sin(45^0+\tfrac{1}{2}S')\cos\tfrac{1}{2}(T'+\lambda)&=\sin(45^0+\tfrac{1}{2}(S-r))\cos\tfrac{1}{2}T\\
\cos(45^0+\tfrac{1}{2}S')\sin\tfrac{1}{2}(T'-\lambda)&=\cos(45^0+\tfrac{1}{2}(S+r))\sin\tfrac{1}{2}T\\
\cos(45^0+\tfrac{1}{2}S')\cos\tfrac{1}{2}(T'-\lambda)&=\cos(45^0+\tfrac{1}{2}(S-r))\cos\tfrac{1}{2}T
\end{aligned}
\]
\indent The third method has some advantage over the other two in the brevity of the calculation, while those two can generally give slightly more accurate results, particularly for \(\lambda\), which they always give with fully sufficient accuracy. The first method, however, determines \(T'\) comparatively imprecisely when it is close to a right angle. If all three results are required with uniform accuracy, sufficient from the point of view of Art. 10,
